\documentclass[a4paper]{article}

\newcommand{\docTitle}{Maschinelles Lernen}

\usepackage{datetime2}  % for dates
\usepackage[utf8]{inputenc} % UTF-8 input encoding
\usepackage[T1]{fontenc} % fixes \textquotedbl with pdflatex/OT1
\usepackage{textcomp}   % additional text symbols for typewriter/upquote
\usepackage{xcolor}     % for gray color
\usepackage{listings}
\usepackage{marginnote} % for horizontal margin notes
\usepackage{parskip}
\usepackage{amsmath,amssymb,mathtools}
\usepackage{everypage}
\usepackage{fancyhdr}
\usepackage{amsfonts}   % Mengen and natural numbers
\usepackage{enumerate}  % Custom Enumerations
\usepackage[inline]{enumitem}   % List spacing and labels
\usepackage{multicol}   % Multiple Columns
\usepackage{tabularx}

\usepackage{hyperref}   % Links
\hypersetup{
    pdftitle={\docTitle},
    pdfauthor={Moritz Köhler},
    pdfkeywords={DHBW, Hochschule},
    colorlinks=true,
    linkcolor=black,
    filecolor=magenta,
    urlcolor=cyan
}

% -------
% Images
% -------

\usepackage{import}
\usepackage{xifthen}
\usepackage{pdfpages}
\usepackage{tikz}
\usetikzlibrary{arrows.meta, positioning}

\newcommand{\incfig}[1]{%
    \def\svgwidth{\columnwidth}
    \import{./fig/}{#1.pdf_tex}
}

% ---------------------------------
% Vertical Side note (Bottom Left)
% ---------------------------------

\newcommand{\verticaldocinfo}{%
  \rotatebox{90}{\tiny\color{gray}\texttt{\DTMtoday\_\docTitle\_\jobname}}%
}

\AddEverypageHook{
  \put(-40pt,-650pt){\verticaldocinfo}%
}

% ---------------
% Math Shortcuts
% ---------------

% Number sets
\newcommand{\R}{\mathbb{R}}
\newcommand{\N}{\mathbb{N}}
\newcommand{\Z}{\mathbb{Z}}
\newcommand{\Q}{\mathbb{Q}}
\newcommand{\C}{\mathbb{C}}

% Vectors and matrices
\newcommand{\vect}[1]{\mathbf{#1}}
\newcommand{\mat}[1]{\mathbf{#1}}

% Derivatives
\newcommand{\dd}{\mathrm{d}}
\newcommand{\dpartial}[2]{\frac{\partial #1}{\partial #2}}
\newcommand{\ddt}{\frac{\dd}{\dd t}}

% Norms and absolute values
\newcommand{\norm}[1]{\left\lVert #1 \right\rVert}
\newcommand{\abs}[1]{\left\lvert #1 \right\rvert}

% Probability & Statistics
\newcommand{\E}{\mathbb{E}}
\newcommand{\Var}{\mathrm{Var}}
\newcommand{\Prob}{\mathbb{P}}

% Fields and Algebras
\newcommand{\K}{\mathbb{K}}
\newcommand{\F}{\mathbb{F}}

% Set Operations
\newcommand{\compl}[1]{#1^{\mathrm{c}}}

% Matrix and Vector Operations
\newcommand{\T}{^{\top}}
\newcommand{\inv}{^{-1}}
\newcommand{\conj}[1]{\overline{#1}}

% -------------------
% Main lecture macro 
% -------------------

\newcommand{\lecture}[1]{%
  \protect\marginnote{\protect\tiny\protect\color{gray}#1}
}

% -----------------------------
% Command for box with heading
% -----------------------------

\fboxsep2mm
\newcommand{\know}[2]{
    \noindent\fbox{
        \parbox{\textwidth-21pt}{

            \textbf{#1:}
            \vspace{0.125cm}

            #2
        }
    }
}

% -------------------
% Listings Languages
% -------------------

% Shared defaults for code listings
\lstset{
    basicstyle=\ttfamily\small,
    keywordstyle=\color{blue}\bfseries,
    commentstyle=\color{gray},
    stringstyle=\color{teal},
    numberstyle=\tiny\color{gray},
    numbers=left,
    stepnumber=1,
    numbersep=8pt,
    showstringspaces=false,
    frame=single,
    breaklines=true,
    tabsize=4,
    keepspaces=false,
    columns=flexible,
    upquote=true,
    extendedchars=false % disable extended char handling to avoid UTF-8 conflicts
}

% SQL syntax highlighting
\lstdefinelanguage{SQL}{
    morekeywords={SELECT,FROM,WHERE,INSERT,INTO,VALUES,UPDATE,SET,DELETE,CREATE,TABLE,PRIMARY,KEY,NOT,NULL,AND,OR,JOIN,LEFT,RIGHT,INNER,ON,AS,ORDER,BY,GROUP,HAVING,LIMIT},
    sensitive=false,
    morecomment=[l]{--},
    morestring=[b]'
}

% JSON syntax highlighting
\lstdefinelanguage{json}{
    morestring=[b]",
    stringstyle=\color{teal},
    comment=[l]{//},
    morecomment=[s]{/*}{*/},
        morekeywords={true,false,null},
    literate=
     *{0}{{{\color{blue}0}}}{1}
      {1}{{{\color{blue}1}}}{1}
      {2}{{{\color{blue}2}}}{1}
      {3}{{{\color{blue}3}}}{1}
      {4}{{{\color{blue}4}}}{1}
      {5}{{{\color{blue}5}}}{1}
      {6}{{{\color{blue}6}}}{1}
      {7}{{{\color{blue}7}}}{1}
      {8}{{{\color{blue}8}}}{1}
      {9}{{{\color{blue}9}}}{1}
            {:}{{{\color{black}:}}}{1}
            {,}{{{\color{black},}}}{1},
    showstringspaces=false
}

% Language specific styles
\lstdefinestyle{javastyle}{
    language=Java,
    basicstyle=\ttfamily\small,
    keywordstyle=\color{blue}\bfseries,
    commentstyle=\color{gray},
    stringstyle=\color{teal},
    numberstyle=\tiny\color{gray},
    numbers=left,
    stepnumber=1,
    numbersep=8pt,
    showstringspaces=false,
    frame=single,
    breaklines=true,
    tabsize=4,
    keepspaces=true,
    columns=fullflexible,
    upquote=true
}

\lstdefinestyle{sqlstyle}{
    language=SQL,
    basicstyle=\ttfamily\small,
    keywordstyle=\color{blue}\bfseries,
    commentstyle=\color{gray},
    stringstyle=\color{teal},
    numberstyle=\tiny\color{gray},
    numbers=left,
    stepnumber=1,
    numbersep=8pt,
    showstringspaces=false,
    frame=single,
    breaklines=true,
    tabsize=4,
    keepspaces=true,
    columns=fullflexible,
    upquote=true
}

\lstdefinestyle{pythonstyle}{
    language=Python,
    basicstyle=\ttfamily\small,
    keywordstyle=\color{blue}\bfseries,
    commentstyle=\color{gray},
    stringstyle=\color{teal},
    numberstyle=\tiny\color{gray},
    numbers=left,
    stepnumber=1,
    numbersep=8pt,
    showstringspaces=false,
    frame=single,
    breaklines=true,
    tabsize=4,
    keepspaces=true,
    columns=fullflexible,
    upquote=true
}

\lstdefinestyle{cstyle}{
    language=C,
    basicstyle=\ttfamily\small,
    keywordstyle=\color{blue}\bfseries,
    commentstyle=\color{gray},
    stringstyle=\color{teal},
    numberstyle=\tiny\color{gray},
    numbers=left,
    stepnumber=1,
    numbersep=8pt,
    showstringspaces=false,
    frame=single,
    breaklines=true,
    tabsize=4,
    keepspaces=true,
    columns=fullflexible,
    upquote=true
}

\lstdefinestyle{bashstyle}{
    language=bash,
    basicstyle=\ttfamily\small,
    keywordstyle=\color{blue}\bfseries,
    commentstyle=\color{gray},
    stringstyle=\color{teal},
    numberstyle=\tiny\color{gray},
    numbers=left,
    stepnumber=1,
    numbersep=8pt,
    showstringspaces=false,
    frame=single,
    breaklines=true,
    tabsize=4,
    keepspaces=true,
    columns=fullflexible,
    upquote=true
}

\lstdefinestyle{jsonstyle}{
    language=json,
    basicstyle=\ttfamily\small,
    keywordstyle=\color{blue}\bfseries,
    commentstyle=\color{gray},
    stringstyle=\color{teal},
    numberstyle=\tiny\color{gray},
    numbers=left,
    stepnumber=1,
    numbersep=8pt,
    showstringspaces=false,
    frame=single,
    breaklines=true,
    tabsize=4,
    keepspaces=true,
    columns=fullflexible,
    upquote=true
}

% Default listing style
\lstset{language={}}

\title{\docTitle}
\author{Moritz}
\date{\today}

\begin{document}
\maketitle
\tableofcontents

\lecture{2026-10-09}

\section{Mathematische Grundlagen}

Es gibt eine Handschriftliche Mitschrift auf Moodle.

\begin{itemize}
    \item Skalar: $x$
    \item Vektoren: $\vect{x}$
    \item Matrizen: $\mat{x}$
\end{itemize}

\subsection{Korrelation und Faltung}

\subsubsection{Korrelation}

% TODO: Formel später mit der Mitschrift abgleichen.

Korrelation

$$R_{xy} = \alpha \sum_{n=1}^{N} x(n) y^*(n)$$

\begin{itemize}
    \item $\alpha$ Konstante $\alpha \neq 1$ oder $\alpha = \frac{1}{N}$
    \item $y^*(n)$ komplexe Konjugation von $y(n)$
\end{itemize}

Kovarianz (Korrelation ohne Mittelwert)

$$C_{xy} = \alpha \sum_{n=1}^{N} (x(n) -\mu_x)(y(n)-\mu_y)$$

Varianz $\sigma^2$

Kovarianzmatrix. Dies sind die Beziehungen zwischen mehreren Variablen (Die Diagonale ist die Variable mit sich selbst.) Optimierung meist so das außerhalb der diagonale die Werte gegen 0 gehen.

Korrelationskoeffizient $\rho$ zwischen $x(n)$ und $y(n)$

$$\rho_{xy} = \frac{C_{xy}}{\sqrt{C_{xx} C_{yy}}} = \frac{C_{xy}}{\sigma_x*\sigma_y}$$

Korreliert nur, wenn es nicht Null ist. Es ist kohärent (Maximale positiv oder negativ) wenn der Betrag = 1 ist. (S.2) Im Beispiel geht es um die Steigung der beiden Signalen.

Korrelationen sind mit Vorsicht zu genießen (Eignet sich zur Auswahl/ Vorverarbeitung)

\begin{itemize}
    \item Der Standard: Pearson-Korrelation
    \item Rangkorrelation (nach Spearman oder Kendall, S. 8) - Hier geht es um den Rang. Ausreißer fallen hier nicht so auf.
\end{itemize}

\know{Korrelation und Kausalität}{Diese sind nicht Gleich. Grafiken S. 10ff oder \href{https://tylervigen.com/spurious-correlations}{Tyler Vigen's project list}, \href{https://web.archive.org/web/20140509212006/http://tylervigen.com/}{Archive}}

\subsubsection{Faltung (Convolution)}

Mathematische Operation auf zwei Funktionen.

$$(x\ast y)(n) = \sum_{m=-\infty}^{\infty}x(m)y(n-m)$$

Vorstellung: Wir invertieren die zweite Funktion und machen einen Moving Window mit Addition aller Produkte.

Die Faltung ist in Bildern (2D Faltung) sehr wichtig. Das sind die Filter wie Kanten oder Schärfe aus Autonomes Fahren.

\subsection{Stochastik}

Ist gleich wie Mathe Semester 4.

\subsubsection{Wahrscheinlichkeiten}

% Ist in der Händischen Mitschrift

Begriffe zum Angucken:

\begin{itemize}
    \item Bedingte Wahrscheinlichkeit
    \item nach Bayes
\end{itemize}

Bayes Theorem:

$$P(A\mid B) = P(B\mid A) * \frac{P(A)}{P(B)}$$

Totale Wahrscheinlichkeit

$$P(A) = P(A\mid B) * P(B) + P(A\mid \overline{B}) * P(\overline{B})$$

Kombination: Bayes-Theorem mit Gesetz der totalen Wahrscheinlichkeit:

% Formel Fehlt mir

\subsubsection{Zufallsvariablen und Wahrscheinlichkeitsverteilungen}

Wir Mappen Zahlen auf Ereignisse.

Begriffe:

\begin{itemize}
    \item Diskrete Zufallsvariable
    \item Kontinuierliche Zufallsvariable
    \item Wahrscheinlichkeitsdichtefunktion/ Probability Density Funktion (PDF)
    \item Verteilungsfunktion
    \item Gaußverteilung/ Normalverteilung
    \item Standardnormalverteilung
    \item Laplaceverteilung
\end{itemize}

\subsection{Grundlagen der Optimierung}

\subsubsection{Was versteht man unter Optimierung?}

Mögliches Optimierungsproblem: Wie kann ich meine Ausgaben so optimieren, dass ich möglichst viel fr meinen Urlaub sparen kann?

Mathematisch: Eine Optimierung ist meist eine Globale Maxi- oder Minimierung die noch zusätzliche Nebenbedingungen erfüllt.

\subsubsection{Beispiel für Optimierungsverfahren}

\begin{itemize}
    \item Grid-Search (Findet Gute Lösung, aber sehr rechenintensiv)
    \item Lineare Optimierung (Lösung von Gleichungssystemen)
    \item Least Square
    \item Gradient Descent
    \item Kreuzentropie-Methode (Ein Monte-Carlo-Verfahren)
    \item Schwarzintelligenz-Methode (Ant-Algorithms)
\end{itemize}

\end{document}